\begingroup
\section{Node DF0B: the comparable-grid discrete square}
\label{module:DF0B}
\noindent\textit{Audited source: \nolinkurl{DF0_good_grid.tex}.}
\subsection{Quasiuniform data and the exact material matrix}

Let $M\ge2$, $b=2\pi/M$, and let
\begin{equation}\label{DF0B:eq:quasi}
 h_i>0,\qquad \sum_i h_i=2\pi,
 \qquad (1-\rho)b\le h_i\le(1+\rho)b,
 \qquad 0<\rho<\frac14.
\end{equation}
Put $d_i=h_i-b$ and $h_i(t)=b+td_i$.  Then
\begin{equation}\label{DF0B:eq:path-quasi}
 (1-\rho)b\le h_i(t)\le(1+\rho)b,
 \qquad |d_i|\le\rho b,\qquad \sum_i d_i=0.
\end{equation}

For a positive current partition $g=(g_i)$ with endpoints $\theta_i$ and
midpoints $c_i$, the assembled foundation F148 gives the exact material
derivative of
\begin{equation}\label{DF0B:eq:w}
 w_i=\Hil q_g(c_i)=\frac1\pi\sum_jg_j\Phi(c_i-c_j),
 \qquad \Phi'(x)=-\log\left|2\sin\frac x2\right|,
\end{equation}
in the form
\begin{equation}\label{DF0B:eq:matrix}
 \dot w_i=\frac1\pi\sum_kd_k\mathcal E_{ik}(g).
\end{equation}
For completeness, choose for each target $i$ a cut $\theta_0$ at circular
distance at least $\pi/3$ from $c_i$.  With
$K=\Phi'$ and cyclic indices relabelled at that cut, the entries are
\begin{align}
 \mathcal E_{ik}^-&=
 \sum_{j<k}g_jK(c_i-c_j)+\frac{g_k}{2}K(c_i-c_k)
 -\int_{\theta_0}^{c_k}K(c_i-y)\,dy, &&k<i,\notag\\
 \mathcal E_{ik}^+&=
 \int_{c_k}^{\theta_0+2\pi}K(c_i-y)\,dy
 -\frac{g_k}{2}K(c_i-c_k)-\sum_{j>k}g_jK(c_i-c_j), &&k>i,
 \label{DF0B:eq:E-explicit}\\
 \mathcal E_{ii}&=\frac12\left\{
 \sum_{j<i}g_jK(c_i-c_j)-\int_{\theta_0}^{\theta_i}K(c_i-y)\,dy
 \right.\notag\\
 &\hspace{34mm}\left.
 -\sum_{j>i}g_jK(c_i-c_j)
 +\int_{\theta_{i+1}}^{\theta_0+2\pi}K(c_i-y)\,dy
 \right\}.\notag
\end{align}
Set $\mathcal E_{ik}=\mathcal E_{ik}^-$ for $k<i$ and
$\mathcal E_{ik}=\mathcal E_{ik}^+$ for $k>i$.  F148 proves
\eqref{DF0B:eq:matrix} by differentiating the Clausen formula: every raw
coefficient equals a common $\Phi(c_i)$ plus the displayed quadrature
primitive, and the common term vanishes because $\sum_kd_k=0$.

Every entry also has the following signed sawtooth representation.  With
\begin{equation}\label{DF0B:eq:eta}
 \eta(y)=\begin{cases}
 -(y-\theta_j),&\theta_j<y<c_j,\\
 \theta_{j+1}-y,&c_j<y<\theta_{j+1},
 \end{cases}
\end{equation}
one has on every complete cell
\begin{equation}\label{DF0B:eq:eta-properties}
 |\eta|\le g_j/2,\qquad \int_{I_j}\eta=0,
\end{equation}
and, on the left or right arc selected by an antipodal cut,
\begin{equation}\label{DF0B:eq:saw}
 \mathcal E_{ik}=\pm\int_{\mathrm{arc}(i,k)}
 \eta(y)K'(c_i-y)\,dy,
 \qquad K'(x)=-\frac12\cot\frac x2.
\end{equation}
The terminal half-cell is part of the integral; it is not separated from
the complete-cell primitive.

\subsection{The variable-grid discrete CZ lemma}

Write $|i-k|_M$ for cyclic index distance.  Since adding a constant to
every entry of one row does not change \eqref{DF0B:eq:matrix} on
$\sum d_k=0$, choose the representative
$\widetilde{\mathcal E}$ whose one-sided primitives are centred at an
antipodal cut.  Thus the two terminal values at the cut are opposite (or
zero when there is a single antipodal index).  We do not subtract a row
average, which could destroy the far-index size.

\begin{lemma}[Signed primitive estimates on a comparable mesh]
\label{DF0B:lem:kernel}
Assume
\begin{equation}\label{DF0B:eq:current-comparable}
 \kappa b\le g_j\le\kappa^{-1}b
\end{equation}
with fixed $0<\kappa<1$.  Then, after choosing the three usual shifted
antipodal cuts, the normalized matrix
$A=b^{-1}\widetilde{\mathcal E}$ satisfies
\begin{align}
 |A_{ik}|&\le\frac{C_\kappa}{1+|i-k|_M},\label{DF0B:eq:size}\\
 |A_{i,k+1}-A_{ik}|+|A_{i+1,k}-A_{ik}|
 &\le\frac{C_\kappa}{(1+|i-k|_M)^2}
 \label{DF0B:eq:smooth}
\end{align}
away from the two harmless cut indices.  At a cut the two one-sided
values obey \eqref{DF0B:eq:size}.  Moreover, for every cyclic index interval
$J$,
\begin{multline}\label{DF0B:eq:testing}
 \sum_{i\in J}\left|\sum_{k\in J}A_{ik}\right|^2
 +\sum_{k\in J}\left|\sum_{i\in J}A_{ik}\right|^2
 \le C_\kappa|J|.
\end{multline}
All constants are independent of $M$ and of cumulative displacement of
the variable grid from the regular grid.
\end{lemma}

\begin{proof}
Quasiuniformity converts physical distance to index distance:
\[
 \operatorname {dist}(c_i,c_k)\asymp_\kappa
 b(1+|i-k|_M)
\]
on the selected arc.  In \eqref{DF0B:eq:saw}, a complete cell has both the
endpoint cancellation and the zero mean in
\eqref{DF0B:eq:eta-properties}.  Subtracting $K'(c_i-c_j)$ on that cell and
using $|(K')'(x)|\le C/\operatorname {dist}(x,2\pi\mathbb Z)^2$
gives
\[
 \left|\int_{I_j}\eta(y)K'(c_i-y)\,dy\right|
 \le \frac{C_\kappa b}{(1+|i-j|_M)^2}.
\]
The terminal half-cell is estimated together with its endpoint atom and
costs $C_\kappa b/(1+|i-k|_M)$.  Summing from the antipode proves
\eqref{DF0B:eq:size}.  Advancing the terminal index by one cell leaves exactly
one complete zero-mean cell, proving the $k$-difference in
\eqref{DF0B:eq:smooth}.  Moving the target by one comparable cell and
subtracting the old kernel value gives the $i$-difference.  The cut
values have physical distance comparable with one and satisfy the same
size estimate directly.

For testing, retain the left/right signs from the antipodally centred
primitive.  Summation by parts with \eqref{DF0B:eq:smooth} shows that a target
at index distance $r$ from the nearer boundary of $J$ contributes at
most
\[
 C_\kappa\left(1+\log\frac{|J|+1}{r+1}\right).
\]
The square of this profile is summable:
\[
 \sum_{r=0}^{|J|-1}
 \left(1+\log\frac{|J|+1}{r+1}\right)^2\le C|J|.
\]
This proves the first test.  For the adjoint test, reverse the source and
target roles in \eqref{DF0B:eq:saw}; the complete-cell zero mean and the
terminal endpoint term are unchanged, so the same argument proves the
second test.  No absolute harmonic row sum is used.
\end{proof}

\begin{lemma}[Finite cyclic discrete $T1$ square]
\label{DF0B:lem:T1}
Any cyclic matrix satisfying \eqref{DF0B:eq:size}--\eqref{DF0B:eq:testing} obeys
\begin{equation}\label{DF0B:eq:T1}
 \|Ax\|_{\ell^2}\le C_\kappa\|x\|_{\ell^2}.
\end{equation}
\end{lemma}

\begin{proof}
Use the three shifted dyadic grids on the cyclic index set and expand
$x$ and a test vector $y$ in normalized Haar functions.  If their
support intervals are disjoint, subtract the kernel at the two interval
centres.  Equation \eqref{DF0B:eq:smooth} gives
\[
 |\langle Ah_I,h_J\rangle|
 \le C_\kappa
 \frac{|I|^{1/2}|J|^{1/2}\min\{|I|,|J|\}}
 {(\operatorname {dist}(I,J)+\max\{|I|,|J|\})^2}.
\]
The two-sided annular sum is Schur summable.  For nested supports, split
the larger Haar function into its constant child value plus a mean-zero
remainder.  The remainder has the same separated estimate; the constant
parts form the two paraproducts.  Their coefficients are Carleson because
\eqref{DF0B:eq:testing} is exactly the square testing condition for $A1_J$
and $A^*1_J$.  The dyadic Carleson embedding theorem bounds both
paraproducts.  The three grids place every pair of cyclic intervals in
one of these separated or nested configurations with comparable parent.
Summing the Haar coefficients proves \eqref{DF0B:eq:T1}.
\end{proof}

\begin{theorem}[Quasiuniform quadrature-primitive square]
\label{DF0B:thm:QPS}
Under \eqref{DF0B:eq:current-comparable}, for every zero-sum direction $d$,
\begin{equation}\label{DF0B:eq:QPS}
 \boxed{
 \sum_i\left|\sum_kd_k\mathcal E_{ik}(g)\right|^2
 \le C_\kappa b^2\sum_kd_k^2.}
\end{equation}
\end{theorem}

\begin{proof}
The chosen row gauge does not change the action on $d$.  Apply
Lemmas~\ref{DF0B:lem:kernel}--\ref{DF0B:lem:T1} to
$A=b^{-1}\widetilde{\mathcal E}$ and rescale.
\end{proof}

\subsection{Endpoint integration inside one equality component}

Along \eqref{DF0B:eq:path-quasi}, Theorem~\ref{DF0B:thm:QPS} and
\eqref{DF0B:eq:matrix} give
\begin{equation}\label{DF0B:eq:material-L2}
 \sum_i|\dot w_i(t)|^2\le C_\rho b^2\sum_kd_k^2.
\end{equation}
At $t=0$ the regular $M$-fan has $w_i(0)=0$.

\begin{theorem}[Complete quasiuniform AVS]
\label{DF0B:thm:AVS}
Every fan satisfying \eqref{DF0B:eq:quasi} obeys
\begin{equation}\label{DF0B:eq:AVS}
 \boxed{
 T(\alpha)=\sum_i h_i|\Hil q_\alpha(c_i)|^2
 \le C_\rho b^3\sum_i(h_i-b)^2
 \le C_\rho\Dvar(\alpha).}
\end{equation}
Hence HQ and DF are closed on the entire quasiuniform component, with no
adjacent-ratio derivative or cumulative-position hypothesis.
\end{theorem}

\begin{proof}
Cauchy--Schwarz in $t$, \eqref{DF0B:eq:quasi}, and
\eqref{DF0B:eq:material-L2} give
\begin{align*}
 T(\alpha)
 &\le(1+\rho)b\int_0^1\sum_i|\dot w_i(t)|^2\,dt\\
 &\le C_\rho b^3\sum_i d_i^2.
\end{align*}
For the last inequality in \eqref{DF0B:eq:AVS}, use the pair formula of stage
153.  On \eqref{DF0B:eq:quasi},
\[
 h_ih_j(h_i^2-h_j^2)^2\asymp_\rho
 b^4(h_i-h_j)^2.
\]
Since $2\pi=Mb$ and $\sum_i d_i=0$,
\[
 \Dvar(\alpha)
 \asymp_\rho\frac{b^4}{Mb}
 \sum_{i,j}(d_i-d_j)^2
 =2b^3\sum_i d_i^2
\]
up to constants depending only on $\rho$.  Stage 153 gives
AVS $\Rightarrow$ HQ $\Rightarrow$ DF.
\end{proof}

The same square is needed at the collapsed locations at which a bad
block will later be opened.  Write $z_i=\theta_i$ for the cell endpoints.

\begin{corollary}[Quasiuniform endpoint samples]
\label{DF0B:cor:endpoints}
Under \eqref{DF0B:eq:quasi},
\begin{equation}\label{DF0B:eq:endpoint-samples}
 \boxed{
 b\sum_i|\Hil q_\alpha(z_i)|^2
 \le C_\rho b^3\sum_i(h_i-b)^2.}
\end{equation}
\end{corollary}

\begin{proof}
Repeat the material calculation with the target carried by the endpoint
$z_i(t)$ instead of the midpoint $c_i(t)$.  The exact length gradient is
again a signed quadrature primitive.  At the target endpoint the two
incident terminal half-cells must be kept together; their quadratic
profiles and first derivatives vanish at their common endpoint.  Thus
there is no boundary logarithm.  Advancing either source or target by one
cell leaves one complete mean-zero sawtooth, so
\eqref{DF0B:eq:size}--\eqref{DF0B:eq:testing} hold for the endpoint matrix with the
same constants.  Lemma~\ref{DF0B:lem:T1} therefore gives the analogue of
\eqref{DF0B:eq:material-L2}.  At the regular fan the spline is even about every
endpoint, so all endpoint Hilbert samples vanish.  Integration in $t$
proves \eqref{DF0B:eq:endpoint-samples}.
\end{proof}

\begin{remark}[Why stage 152 does not contradict this theorem]
The false HMS/QPS family has half of its current cells of length
$\varepsilon a$ while their fixed straight-path directions have size
$a$.  It leaves every fixed quasiuniform component.  In
Theorem~\ref{DF0B:thm:AVS}, both the current lengths and the entire connecting
segment remain comparable with $b$, and $|d_i|\le\rho b$.
\end{remark}

\subsection{Effect on the ELCA ledger}

Use the good/bad selection of stage 153.  On a maximal collection of
good cells whose zero-cell deletion contains no intervening bad block,
the current effective mesh is comparable with the selected stratum
length $b$.  The preceding theorem closes every good--good
first-separation rectangle.

\begin{corollary}[Only bad-block rectangles remain]
\label{DF0B:cor:remaining}
In ELCA of stage 153, all rows whose source and target shadows avoid the
bad set satisfy
\begin{equation}\label{DF0B:eq:good-good}
 T_{\mathrm{good\text{-}good}}
 \le C_\rho b^3\sum_{i\in\mathcal G_\rho}(h_i-b)^2.
\end{equation}
Thus the remaining pure endpoint ledger consists only of rectangles for
which the source shadow or target shadow meets a maximal bad block.  Such
a block must be retained as a whole until its first separating scale and
paid by
\begin{equation}\label{DF0B:eq:bad-charge}
 \sum_{i\in\mathcal B_\rho}h_i(h_i^2-\ell^2)^2.
\end{equation}
\end{corollary}

\begin{proof}
Apply Theorem~\ref{DF0B:thm:AVS} on each good component after the common
collapsed-uniform reference has been fixed.  At its two boundary cuts,
keep the terminal half-cell with the adjacent bad block; those boundary
rectangles are not included in \eqref{DF0B:eq:good-good}.  Disjoint good
components have disjoint physical mass, so their squares sum.  The
reference-charge estimate of stage 153 then gives the displayed right
side.  Every omitted rectangle meets a bad block by construction.
\end{proof}
\endgroup
